
\subsection{Problem Statement}
Automated Speaker Identification (SID) plays a pivotal role in modern biometric security and human-computer interaction. However, the reliability of such systems is often compromised by the inherent variability of human speech and the presence of environmental noise . Speech is a complex, non-stationary signal that conveys both linguistic content and speaker-specific physiological traits. Identifying a speaker accurately requires a system that can distinguish these unique vocal characteristics from background interference and channel distortions .
\subsection{Motivation}
Although modern Deep Learning (DL) models have shown remarkable capability to learn features directly from data, they often require massive datasets and high computational power to achieve robustness. Integrating Digital Signal Processing (DSP) into the AI pipeline offers a more efficient path by "cleaning" the signal and extracting handcrafted features that represent the physical properties of the human vocal tract. This hybrid approach—combining classical DSP with ML—is motivated by the need for systems that are not only accurate, but also computationally sustainable and interpretable.
\subsection{Research Question}
The central inquiry of this report is: To what extent does a dedicated DSP preprocessing pipeline (Pipeline B) improve the classification accuracy, Signal-to-Noise Ratio (SNR), and generalization of AI models compared to a baseline approach using raw signal inputs (Pipeline A)?